% !TEX root = ../main.tex
% !TEX program = lualatex
\documentclass[../main.tex]{subfiles}
\IfSubfilesClassLoaded{\externaldocument{\subfix{../build/main}}}{}
% =============================================================================
% File: 27F_MBTI.tex
% =============================================================================
\begin{document}
\IfSubfilesClassLoaded{\chapter{Myers-Briggs Type Indicator}}{}
\section{Myers-Briggs Type Indicator}

\subsection{Summary}
\begin{description}
  \item[\textbf{Preference framework.}] The \ac{mbti} describes preferences in how people energize, gather information, decide, and adapt to the outer world; it is a sorting tool, not a measure of skill or ability~\autocite{edo-5-1-8-mbti-2025}.
  \item[\textbf{Self-awareness.}] The goal is non-judgmental feedback that improves self-awareness and helps leaders accommodate different work-style preferences~\autocite{edo-5-1-8-mbti-2025}.
  \item[\textbf{Team impact.}] Similar types communicate quickly but can miss perspectives; diverse types reduce error risk when leaders build understanding and trust~\autocite{edo-5-1-8-mbti-2025}.
\end{description}

\subsection{History and Purpose}
\begin{description}
  \item[\textbf{Origins.}] Carl Jung described predictable patterns in how people prefer to engage the world; Isabel Myers and Katherine Briggs translated those preferences into the \ac{mbti} to help people match skills, desires, and career opportunities~\autocite{edo-5-1-8-mbti-2025}.
  \item[\textbf{Use case.}] The \ac{mbti} sorts preferences to support self-awareness and leadership effectiveness, not to evaluate performance or select personnel~\autocite{edo-5-1-8-mbti-2025}.
\end{description}

\subsection{The MBTI Dichotomies}
\begin{description}
  \item[\textbf{Energizing.}] Extraversion (outer world) vs. introversion (inner world) preferences define where people direct and receive energy~\autocite{edo-5-1-8-mbti-2025}.
  \item[\textbf{Collecting data.}] Sensing focuses on facts and experience; intuition focuses on patterns and possibilities~\autocite{edo-5-1-8-mbti-2025}.
  \item[\textbf{Evaluating.}] Thinking emphasizes logic and objectivity; feeling emphasizes values and impact on people~\autocite{edo-5-1-8-mbti-2025}.
  \item[\textbf{Adapting.}] Judging favors structure and closure; perceiving favors flexibility and openness to options~\autocite{edo-5-1-8-mbti-2025}.
\end{description}

\subsection{Self-Select and Best-Fit Types}
\begin{description}
  \item[\textbf{Self-select vs. reported.}] Compare self-estimated type with reported \ac{mbti} results; if different, review both descriptions and decide which fits best~\autocite{edo-5-1-8-mbti-2025}.
  \item[\textbf{Best-fit hierarchy.}] Best-fit type is the goal; reported type and self-estimated type inform the decision, but you decide what fits~\autocite{edo-5-1-8-mbti-2025}.
  \item[\textbf{Why differences occur.}] Mindset, changes over time, and familiarity with questions can shift responses; differences are not a cause for concern~\autocite{edo-5-1-8-mbti-2025}.
\end{description}

\subsection{MBTI Step II Insights}
\begin{description}
  \item[\textbf{Preference confidence.}] The preference confidence index and category indicate how consistently a person selected one side of a dichotomy; they suggest, but do not determine, accuracy~\autocite{edo-5-1-8-mbti-2025}.
\end{description}

\subsection{Team Dynamics and Problem Solving}
\begin{description}
  \item[\textbf{Team composition.}] High similarity can accelerate understanding but risks incomplete viewpoints; diverse types can reduce conflict once members learn to value differences~\autocite{edo-5-1-8-mbti-2025}.
  \item[\textbf{Problem-solving flow.}] Effective teams balance sensing and intuition for data collection, and thinking and feeling for evaluation to avoid blind spots~\autocite{edo-5-1-8-mbti-2025}.
  \item[\textbf{Self-discovery.}] The Johari Window reminds teams that feedback and transparency reduce blind spots and build trust~\autocite{edo-5-1-8-mbti-2025}.
\end{description}

\subsection{Quick Review}
\begin{itemize}
  \item \ac{mbti} describes preferences, not competence or ability~\autocite{edo-5-1-8-mbti-2025}.
  \item Leaders should adapt communication to different preferences to improve performance~\autocite{edo-5-1-8-mbti-2025}.
  \item Diverse types improve decision quality when teams value differences and share information~\autocite{edo-5-1-8-mbti-2025}.
\end{itemize}

\IfSubfilesClassLoaded{
  \RenewDocumentCommand{\entryname}{}{\textbf{\color{Modern} Acronym}}
  \RenewDocumentCommand{\descriptionname}{}{\textbf{\color{Modern} Definition}}
  \printnoidxglossary[
    type=\acronymtype,
    title=Acronyms,
    style=long-booktabs
  ]}{}
\end{document}
